\paragraph{Asaeda--Haagerup and recent odd-order constructions.}
The ring \sha{R8_3_721} has the exact Asaeda--Haagerup global dimension
\(136+32\sqrt{17}\), so a dimension mismatch does not dispose of this
lead.  The complete known Morita class nevertheless gives no model:
\(AH_1,AH_2,AH_3\) have ranks \(6,9,9\), respectively
\cite[Section~4.3, Tables~1--3]{GSAsaedaMorita}, while
\(AH_4,AH_5,AH_6\) have four invertibles and maximum multiplicity two
\cite[Lemmas~3.4, 3.7, 3.15 and Theorem~4.21]{AH}.
The source ring has one invertible and maximum multiplicity three.
The audit independently reconstructs the two rank-eight AH based
tensors and matches all \(512\) coefficients to the existing positive
controls \sha{R8_2_39} and \sha{R8_2_484}.  Of the \(277\) baseline
survivors, fourteen have the AH global dimension, checked using exact
positive dimension characters; none is a member of this known family.

The September 2026 constructions
\cite[Theorem~2.1]{HuangOddHI2026} and
\cite[Theorems~A, 4.4 and 6.3]{GSYOddHI2026} were also checked.
For odd \(N\geq3\), write \(d^2=Nd+1\) with \(d>0\).
The HI family has global dimension \(D_N=N(1+d^2)\), and the specified
near-group family has global dimension \(ND_N\).
Exact Casimir characteristic-polynomial division, supplemented by
integer spectral bounds, rules out both dimensions for every baseline
survivor and every odd \(N\); the bounds \(D_N>N^3\) and
\(ND_N>N^4\) certify the infinite tail.  This also rules out their
Morita duals, since Morita equivalence preserves FP dimension.
The cited general existence statements are spherical and pseudo-unitary;
the stated Morita classification has additional unitary, cyclic,
untwisted and Q-system hypotheses.  These completed model-family
checks make no general nonexistence claim for the remaining rings.
